\scanpage{52}
\mspage{25}
7. \uline{Étude des} \deleted{\unsure{cas}} $\widehat M_G$ \deleted{\unsure{en dim}} \uline{de dim} $\leq 2$.

a) Détermination des MD-graphes correspondants.

I) \uline{Cas de la dimension $0$.} Les \deleted{MD-\unsure{gr}} \inserted{conditions} suivantes sur le MD-graphe $\underline G$ sont équivalentes

$a)$ $\dim M_G=0\qquad(\Longleftrightarrow\dim\widehat M_G=0)$

$b)$ \deleted{Et} la ``dim. modulaire'' de $\underline G$ est nulle, i.e. le nb $s_1$ d'arêtes de \deleted{$\unsure{\underline G}$} $G$ (la codimension modulaire de $\underline G$) est égale à son maximum $3(g-1)+\nu$.

$c)$ \deleted{\unsure{Toute}\ \illegible{1}} $\underline G$ est minimal parmi les MD-graphes pour la relation de générisation, i.e. tt morphisme de générisation $\underline G'\mathbin{\text{---}}\underline G$ de MD-graphes est un iso.

$d)$ Tout sommet \inserted{$\alpha$} de $\underline G$ est de genre $0$, et est d'ordre \deleted{$\unsure{3}$} $\omega_\alpha\leq3$; \deleted{\unsure{genre}\ \unsure{le poids}} \inserted{le poids marqué} $\nu_\alpha$ est égal à $3-\omega_\alpha$.

On peut alors, \uline{à isom unique près}, trouver une courbe \unsure{standart|standard} \inserted{$\underline G$-épinglée} et une seule sur une base quelconque $S$, \deleted{$\underline G$-épinglée}, i.e.
\[
\text{(60)}\qquad M_G=\unsure{M_{\widehat G}|\widehat M_G}\simeq\operatorname{Spec}\mathbf Z\overset{\mathrm{def}}=S_0,
\]
donc
\[
\text{(61)}\qquad M_{[G]}\simeq B(\Gamma_{S_0})
\]
\margin{(\deleted{\unsure{espace}} \inserted{multiplicité} classifiante pour le schéma en groupes $\Gamma_{S_0}$ sur $S_0=\operatorname{Spec}\mathbf Z$, où $\Gamma=\operatorname{Aut}\underline G$)}
\margin{51}
